\documentclass[11pt,a4paper,sans]{moderncv}

% moderncv themes
\moderncvtheme[grey]{casual}                 % optional argument are 'blue' (default), 'orange', 'red', 'green', 'grey' and 'roman' (for roman fonts, instead of sans serif fonts)
%\moderncvtheme[green]{casual}                % idem

% character encoding
\usepackage[english]{babel}
\usepackage[utf8]{inputenc}                   % replace by the encoding you are using

\usepackage{lmodern}
%\renewcommand{\sfdefault}{\rmdefault}
% adjust the page margins
\usepackage[scale=0.75]{geometry}
\usepackage{graphicx}
%\setlength{\hintscolumnwidth}{3cm}						% if you want to change the width of the column with the dates
%\AtBeginDocument{\setlength{\maketitlenamewidth}{6cm}}  % only for the classic theme, if you want to change the width of your name placeholder (to leave more space for your address details
%\AtBeginDocument{\recomputelengths}                     % required when changes are made to page layout lengths

%----------------------------------------------------------------------------------
%            Kontaktdaten
%----------------------------------------------------------------------------------
% VORNAME
\firstname{lukas}
% NACHNAME
\familyname{elsner}
%TITEL (optional, ggf. einfach die Zeile löschen!)
\title{curriculum vitae}
%ADRESSE  (optional, ggf. einfach die Zeile löschen!)
%\address{dorfstrasse 11}{25485 langeln}{germany}
%HANDYNUMMER  (optional, ggf. einfach die Zeile löschen!)
%\mobile{0176 20 21 52 67}
%\mobile{(+49) 176 20 21 52 67}
%EMAIL-ADRESSE  (optional, ggf. einfach die Zeile löschen!)
\email{open@mindrunner.de - cv@lukaselsner.de}
%HOMEPAGE  (optional, ggf. einfach die Zeile löschen!)
\homepage{www.mindrunner.de - www.lukaselsner.de}
%FOTO  (optional, ggf. einfach die Zeile löschen!)
%  64pt = Höhe des Bildes, 'picture' = Name des Bildes
\photo[128pt]{picturesmaller}

% to show numerical labels in the bibliography; only useful if you make citations in your resume
\makeatletter
\renewcommand{\bibliographyitemlabel}{\@biblabel{\arabic{enumiv}}}
\makeatother

%----------------------------------------------------------------------------------
%            Inhalt
%----------------------------------------------------------------------------------

\begin{document}
\maketitle
\section{}
\cvline{born}{01.11.1984 in Henstedt-Ulzburg, Germany}

\section{education}
\cventry{2014--2015}{Master of Information Technology}{Queensland University of Technology}{Brisbane}{Australia}{
  \begin{itemize}
      \item Graded: 7 (High Distinction)
    \item Software Architecture (Advanced)
  \end{itemize}
}
\cventry{2015}{Master Thesis}{\textbf{[Pyhton, C, Java, VHDL, git, \LaTeX]}}{QUT, Schatz Forensic}{}{
  \begin{itemize}
    \item In this project, a low level framework with the goal of reading raw
    NAND flash devices found on SD-Cards and USB-drives was created. The
    hardware/software co-design runs on a Pipistrello FPGA
    board and enables dynamic pin assignment, high speed block
    orientated I/O and voltage variation to particular pins on the
    hardware side, as well as a serial interface to an attached remote machine
    for configuring the hardware and read data from the NAND flash on the
    software side.
  \end{itemize}
}
\cventry{2013--2014}{M.Sc. Information Technology}{University}{Lübeck}{Master of Science}{
  \begin{itemize}
    \item Software Systems Engineering
  \end{itemize}
}
\cventry{2013}{Software Systems Engineering Laboratory}{\textbf{[Java, Maven, Jenkins, git, \LaTeX]}}{Uni Lübeck}{}{
  \begin{itemize}
    \item required course for every Software Systems Engineering student
  \end{itemize}
}
\cventry{2009--2012}{B.Sc. Information Technology}{University}{Passau}{Bachelor of Science}{
  \begin{itemize}
    \item Graded 1.9
    \item Minor: Business Administration
  \end{itemize}
}
\cventry{2012}{Bachelor Thesis}{\textbf{[C++, \LaTeX]}}{Uni Passau}{\url{http://www.uni-passau.de}}{
  \begin{itemize}
    \item planning and implementing an easy-to-use class library on top of a proprietary SDK for prototyping perception platforms with the goal to automate testing within the whole research project (Graded: 1.3)
  \end{itemize}
}
\cventry{2012}{Research Assistant}{\textbf{[C++]}}{Uni Passau}{}{
  \begin{itemize}
    \item developing and troubleshooting of plugins for a perception platform
    \item developing of algorithms for automated image and pattern recognition
  \end{itemize}
}
\cventry{2011}{Software Engineering Laboratory}{\textbf{[Java, Maven, Jenkins, git, \LaTeX]}}{Uni Passau}{}{
  \begin{itemize}
    \item Graded 1.0
    \item required course for every Information Technology student
    \item funCKit is an editor and simulator for digital circuits
    \item the source is released under the GPL\\ (\url{https://github.com/mindrunner/funCKit})
  \end{itemize}
}
\cventry{2005--2008}{Theoretical Vocational Training}{Vocational School}{Elmshorn}{IT specialist}{
  \begin{itemize}
    \item Software Developer
  \end{itemize}
}
\cventry{2005--2008}{Practical Vocational Training}{TecBOS GmbH}{Halstenbek}{IT specialist}{
  \begin{itemize}
    \item Software Developer
  \end{itemize}
}
\cventry{1995--2004}{High School}{Dietrich-Bonhoeffer-Gymnasium}{Quickborn}{}{}

\section{work experience}
\cventry{since 2004}{Software Engineer / IT-Consultant}{mindrunner}{}{}{
  \begin{itemize}
    \item conducting freelance work as a software engineer and consultant
  \end{itemize}
}
\cventry{2010--2011}{Software Engineer}{freifalt GmbH}{Passau/Hamburg}{}{
  \begin{itemize}
    \item design and implementation of iOS-Apps and their backends
    \item responsible for the company's internal infrastructure, a virtualized, heterogeneous network
  \end{itemize}
}
\cventry{2005--2009}{Software Engineer}{MSA Auer GmbH}{Halstenbek/Berlin}{}{
  \begin{itemize}
    \item developing several core products of the company
  \end{itemize}
}

\section{projects}
\cventry{2015}{Android Application}{\textbf{[Java, Android, Gradle, git-flow, Jira]}}{}{}{
  \begin{itemize}
    \item development of an android app for a social network startup company
    \item focus on cutting edge android technology
    \item REST client / ContentProvider / SyncProvider / material design
  \end{itemize}
}
\cventry{2014}{TouchIt}{\textbf{[Java, Android, Gradle, git, libGDX]}}{}{}{
  \begin{itemize}
    \item development of an android game based on libGDX
    \item published in Google's play store
    \item published the source code (GPLv3)
    \item \url{https://github.com/mindrunner/touch-it}
  \end{itemize}
}
\cventry{2013--2014}{brightup GmbH}{\textbf{[Java, C, .NET, OpenEmbedded, git-flow, Jira, Jenkins, BuildBot, Android, iOS]}}{}{http://www.brightup.de}{
  \begin{itemize}
    \item development of hardware and software for brightup’s central unit and z-wave plugs
    \item maintenance of general development tools and work flows (bug tracker, version control system, continuous integration)
    \item consulting for cloud, smartphone applications and general technical decisions
  \end{itemize}
}
\cventry{2013}{PhotoBooth and Quickshot}{\textbf{[Java, Redmine, Jenkins, EDSDK (Canon)]}}{}{}{
  \begin{itemize}
    \item two kiosk machine applications for a photo studio
  \end{itemize}
}
\cventry{2013}{hantPlayer}{\textbf{[Objective-C, Redmine, AppKit]}}{}{\url{http://www.hantmade.com}}{
  \begin{itemize}
    \item a native video player with special tagging functions
  \end{itemize}
}
\cventry{2013}{Stomt}{\textbf{[Java, MongoDB, Spring, git-flow]}}{}{\url{http://www.stomt.com}}{
  \begin{itemize}
    \item consulting and the developing of backend and Android-App via pro bono
  \end{itemize}
}
\cventry{2012}{HydrostopEU}{\textbf{[Joomla, PHP]}}{}{\url{http://www.hydrostopeu.com}}{
  \begin{itemize}
    \item complete webpage redesign
  \end{itemize}
}
\cventry{2012}{Nordflair Ticketsystem}{\textbf{[Java, Maven, Jenkins, JPA]}}{}{\url{http://nordflair-events.de/}}{
  \begin{itemize}
    \item co-working on a web-based ticket and entry management system
    \item developing the Java-Frontend entry terminal
  \end{itemize}
}
\cventry{2012}{Templogger}{\textbf{[Java, Maven, Jenkins, JPA]}}{}{}{
  \begin{itemize}
    \item custom service which receives data from GSM-Datalogger: once received, it stores the information into a database and sends warning e-mails according to threshold values
  \end{itemize}
}
\cventry{2010--2011}{usAR-Hamburg}{\textbf{[Java, Maven, Jenkins, JPA, iOS]}}{freifalt GmbH}{\url{http://www.freifalt.com}}{
  \begin{itemize}
    \item iOS-App which shows up public transportation based on augmented reality
  \end{itemize}
}
\cventry{2005--2009}{TecBOS.replication}{\textbf{[Delphi]}}{TecBOS GmbH}{\url{http://tecbos.com}}{
  \begin{itemize}
    \item developing a set of windows services with the aim to automatically replicate multiple TecBOS-Databases
  \end{itemize}
}
\cventry{2005--2009}{TecBOS.mobile}{\textbf{[C, Sqlite]}}{TecBOS GmbH}{\url{http://tecbos.com}}{
  \begin{itemize}
    \item Mobile Windows CE Devices can exchange data with the TecBOS.solutions
  \end{itemize}
}
\cventry{2005--2009}{TecBOS.webaccess}{\textbf{[C\#, firebird]}}{TecBOS GmbH}{\url{http://tecbos.com}}{
  \begin{itemize}
    \item porting the native TecBOS.solutions into a web application
  \end{itemize}
}
\cventry{2005--2009}{TecBOS.solutions}{\textbf{[Delphi, firebird]}}{TecBOS GmbH}{\url{http://tecbos.com}}{
  \begin{itemize}
    \item maintaining and developing parts of the application as part of my apprenticeship and subsequent employment
  \end{itemize}
}
\cventry{2008}{CMS}{\textbf{[MySQL, PHP]}}{}{}{
  \begin{itemize}
    \item developing a whole content management system within one week
  \end{itemize}
}
\cventry{2005--2008}{Web shop}{\textbf{[MySQL, PHP]}}{fontinform GmbH}{\url{http://www.fontinform.com}}{
  \begin{itemize}
    \item planning and development of a web shop for digital fonts
  \end{itemize}
}
\section{tech skills}
\includegraphics[width=0.7\textwidth]{skillcloud}
%\cvline{Basic}{GTK, Perl, Python, VHDL}{}{}
%\cvline{Advanced}{Android, Bash, C, C++, C\#, Delphi, iOS, Maven, MongoDB, MySQL, PHP}{}{}
%\cvline{Expert}{Debian, Gentoo, Git, OpenEmbedded, Java, \LaTeX, Proxmox, Ubuntu, Yocto}{}{}
%\cvline{Tools}{CMake, Eclipse, Intellij IDE, Make, Netbeans, Vim}{}{}
%\cvline{Miscellaneous}{Office, OS X, Windows}{}{}

\section{languages}
\cvline{German}{Native}{}
\cvline{English}{Fluent}{}
\cvline{French}{Basic}{}

\section{personal interests}
\cvline{}{\begin{itemize}
    \item climbing
    \item bouldering
    \item running
    \item mountain biking
    \item travelling
\end{itemize}}
~\\
~\\
~\\
Langeln, \today\\ % Aktuelles Datum und Stadt
\end{document}

